For a prime p, define
  f(p) := the least integer n >= 1 such that p divides (n! + 1).
Wilson's theorem implies (p-1)! ≡ -1 (mod p), hence p divides ((p-1)!+1), so f(p) is always defined and f(p) <= p-1.

Question A: Are there infinitely many primes p with f(p) = p-1?
Question B: Does f(p)/p -> 0 for almost all primes p?
(Here "almost all primes" means: #{p <= x : the property fails} = o(pi(x)) as x -> infinity.)
